\section*{Abstract}
\addcontentsline{toc}{section}{Abstract}

Surgical triplet segmentation requires both the spatial localisation of surgical instrument instances and the identification of the \ivttriplet interaction associated with each instance. Chapter~\ref{chap:triplet-grounding} demonstrated that strong instance supervision substantially improves grounding, but also showed that semantic interaction prediction remains difficult. This chapter therefore separates these capabilities. A specialist Mask2Former first grounds instrument instances, after which a multimodal large language model predicts a verb and anatomical target for every grounded instance. Direct supervised fine-tuning (Direct SFT) establishes a viable second-stage baseline. Under ground-truth grounding, Sequential Dialogue Prediction and three-frame temporal context produce promising observed improvements in complete-triplet AP, whereas five-frame context, prompt-level anatomical priors, ontology guidance, and distilled-reasoning supervision do not improve performance consistently. In the complete pipeline, errors in predicted grounding reduce performance, but the two-stage Direct SFT system obtains substantially stronger component-level scores and comparable complete-triplet AP relative to TargetFusionNet. These results establish two-stage multimodal interaction prediction as a viable formulation while showing that single-seed evaluation, optimisation sensitivity, and limited predicted-grounding experiments preclude definitive claims of superiority.

\section{Introduction}
\label{sec:c5:introduction}

The preceding chapters develop a progression from broad surgical scene understanding to spatially grounded instrument--tissue interactions. Chapter~\ref{chap:background} situates action understanding within multitask surgical vision; Chapter~\ref{chap:cholecinstanceseg} provides the instrument-instance annotations required to distinguish individual tools; and Chapter~\ref{chap:triplet-grounding} combines those spatial annotations with surgical action triplets to establish CholecTriplet-Seg and TargetFusionNet. Together, these contributions show that accurate instrument grounding is achievable, but that identifying the verb and, especially, the anatomical target associated with each instrument remains substantially more difficult.

This imbalance motivates a two-stage formulation. Instrument instance segmentation is a dense prediction problem for which specialist architectures are well suited. Once the instruments have been localised, interaction prediction is instead a structured semantic task: the model must interpret the action, anatomy, and surrounding scene for each known instance. Requiring one architecture to learn both capabilities jointly may not be necessary. Separating them also provides a controlled setting in which interaction prediction can first be studied with ground-truth instrument instances and then evaluated as part of a complete pipeline with predicted instances.

Multimodal large language models (MLLMs) offer a natural candidate for the second stage because they combine visual input with textual instructions and can generate structured outputs. This chapter investigates whether an MLLM can serve as a viable interaction predictor when instrument instances are supplied, rather than asking it to perform dense localisation. The resulting framework uses a fixed Mask2Former for instrument grounding and Qwen3-VL-8B-Instruct~\cite{bai2025qwen3vl} for verb--target prediction. The complete triplet is formed by combining the fixed instrument identity from Stage~1 with the interaction predicted by Stage~2.

The study first establishes Direct SFT as the principal baseline and then evaluates controlled changes to the second-stage formulation. Sequential Dialogue Prediction decomposes verb and target prediction into two dialogue turns; temporal variants supply short windows of neighbouring frames; an anatomical-context variant supplies coarse EndoViT-derived evidence around an estimated instrument tip; ontology-guided prompting supplies valid verb--target combinations; and distilled-reasoning supervision uses rationales generated by a larger teacher model. These variants are evaluated as individual adaptations, not as a factorial or cumulative ablation study.

The chapter addresses four research questions:
\begin{enumerate}
    \item Can an MLLM serve as a viable second-stage interaction predictor when instrument instances are provided?
    \item How do sequential factorisation, short temporal context, anatomical priors, ontology guidance, and distilled-reasoning supervision affect interaction prediction?
    \item Can the complete two-stage system remain competitive when the instrument instances are predicted?
    \item Which limitations prevent the observed differences from being interpreted as definitive improvements?
\end{enumerate}

The principal contribution is the two-stage formulation and its empirical evaluation. The results show that Direct SFT is a strong and difficult-to-improve baseline; that Sequential Dialogue Prediction and three-frame context are promising for complete-triplet prediction under ground-truth grounding; and that additional context or longer generated supervision does not automatically improve performance. The chapter further demonstrates that stronger component metrics do not necessarily yield a commensurate increase in complete-triplet AP because the instrument, verb, and target must all be correct for the same instance.

\section{Related Work}
\label{sec:c5:related-work}

\subsection{Grounded Surgical Triplet Prediction}

Surgical action triplets represent an interaction as an \ivttriplet tuple. TripNet introduced triplet recognition as a frame-level task, while Rendezvous subsequently modelled relationships between instruments, verbs, and targets through attention mechanisms~\cite{1nywoye_rec_action_triplets,NWOYE2022102433_rendevous}. Later work investigated temporal modelling, self-distillation, and multi-task optimisation, including MT4MTL-KD and TERL~\cite{gui2023mt4mtl,gui2024tail}. These methods improved frame-level recognition but did not directly associate every interaction with a precise instrument instance.

Chapter~\ref{chap:triplet-grounding} addresses this limitation by introducing CholecTriplet-Seg and triplet segmentation~\cite{alabi2026grounding}. Its Mask2Former-Triplet baseline and TargetFusionNet jointly learn instrument localisation and semantic triplet labels, with TargetFusionNet incorporating weak anatomical priors to support target prediction. The present work retains the same grounded task and evaluation protocol but decomposes localisation and interaction prediction into specialist stages.

\subsection{Multimodal Models in Surgical Vision}

MLLMs have increasingly been studied for surgical visual question answering, dialogue, and scene interpretation. SurgicalGPT combines visual representations with an autoregressive language model for surgical visual question answering, while Surgical-LLaVA adapts instruction-tuned vision--language models to surgical images and videos~\cite{seenivasan2023surgicalgpt,jin2024surgical}. Related triplet-recognition work has also used MLLM-generated semantic descriptions as auxiliary information within discriminative models~\cite{zhang2025mejo}. In contrast, this chapter uses an MLLM itself as the structured interaction-prediction stage after specialist instrument grounding.

\section{Two-Stage Multimodal Interaction Prediction}
\label{sec:c5:method}

\subsection{Problem Formulation}
\label{sec:c5:problem}

Let $I\in\mathbb{R}^{H\times W\times3}$ denote a laparoscopic image. Stage~1 produces a set of grounded instrument instances
\begin{equation}
    \mathcal{S}=\{s_i\}_{i=1}^{N}, \qquad s_i=(m_i,c_i,\mathrm{id}_i),
    \label{eq:c5:grounded-instances}
\end{equation}
where $m_i$ is the binary instrument mask, $c_i$ is its fixed instrument category, and $\mathrm{id}_i$ is an image-level identifier. Stage~2 predicts
\begin{equation}
    f_{\theta}(I,\mathcal{S})\rightarrow\mathcal{Y}, \qquad
    \mathcal{Y}=\{(v_i,t_i)\}_{i=1}^{N},
    \label{eq:c5:interaction-mapping}
\end{equation}
where $v_i$ and $t_i$ denote the verb and target. The final triplet for instance $i$ is $\tau_i=(c_i,v_i,t_i)$. Stage~2 receives the instrument identity as a fixed input and cannot revise it.

Two grounding settings are used. Ground-truth grounding supplies the annotated masks, categories, and instance count, isolating Stage~2 from localisation errors. Predicted grounding supplies the outputs of Stage~1 and therefore evaluates the complete system. False-positive Stage~1 instances are passed to Stage~2, whereas missed instruments are absent and cannot be recovered.

\begin{figure}[htb!]
    \centering
    \includegraphics[width=0.95\linewidth]{figures/two_stage_pipeline.pdf}
    \caption{The two-stage framework. Stage~1 uses Mask2Former to localise and classify instrument instances. Stage~2 receives a single image containing translucent instance overlays, bounding boxes, and textual identifiers, and uses Qwen3-VL-8B-Instruct to predict one verb--target interaction per grounded instance. The complete triplets combine the fixed Stage~1 instrument classes with the Stage~2 predictions.}
    \label{fig:c5:framework}
\end{figure}

\subsection{Stage 1: Fixed Instrument Grounding}

Stage~1 uses the Mask2Former architecture, backbone, augmentation, optimisation, checkpoint-selection criterion, and inference configuration described for instrument instance segmentation in Chapter~\ref{chap:cholecinstanceseg}. For the present study, it is retrained on the CholecTriplet-Seg training split, including the corrected instrument masks available for this task-specific subset. The detector is fixed throughout all Chapter~5 experiments. Predicted masks are retained using the established Mask2Former confidence filtering and are passed to Stage~2 without refinement. Ground-truth and predicted masks use the same overlay-generation procedure; only their source differs.

\subsection{Stage 2: Structured Interaction Prediction}

For single-frame experiments, the Stage~2 input is one RGB image onto which every grounded instrument mask is rendered translucently. Each instance additionally receives a bounding box and an identifier comprising its instrument class and an image-level index, such as \texttt{hook\_1} or \texttt{grasper\_2}. Instrument classes have fixed colours; multiple instances of the same class share a colour and are distinguished by their identifiers. Indices follow the order in which instances are received from the annotation or detector and have no semantic meaning.

The textual instruction lists the allowable verb and target labels and specifies a structured response containing an identifier, verb, and target for every supplied instance. The label vocabulary includes the \texttt{null\_verb} and \texttt{null\_target} categories introduced in Chapter~\ref{chap:triplet-grounding}. All instances in an image are processed jointly in one autoregressive response. Generation is unconstrained: Direct SFT does not enforce valid triplets during decoding or apply ontology-based filtering. Malformed outputs, invalid labels, duplicate identifiers, and omitted instances are treated as missing predictions and receive no credit.

\subsection{Direct Supervised Fine-Tuning}
\label{sec:c5:direct-sft}

Direct SFT trains Qwen3-VL-8B-Instruct to generate the structured verb--target response directly from the grounded visual input and instruction. The base model is loaded with 4-bit NF4 quantisation and bfloat16 computation. QLoRA adapters are applied to the attention and MLP projection layers, with rank $r=64$, scaling factor $\alpha=128$, and dropout $0.05$; the vision encoder, multimodal projector, and base language model remain frozen. Supervised loss is computed only over assistant-response tokens.

Training uses AdamW with an initial learning rate of $2\times10^{-4}$, weight decay $0.01$, a warm-up ratio of $0.03$, and linear decay to zero. Gradients are clipped to a maximum norm of $1.0$. A per-device batch size of one and gradient accumulation over 32 steps give an effective batch size of 32. Models are trained for at most five epochs using seed 42, and the checkpoint with the lowest validation loss is reported. All non-temporal experiments use the same visual-token budget. Evaluation uses the same standard, unconstrained autoregressive generation procedure for every adaptation unless stated otherwise.

\subsection{Interaction-Prediction Adaptations}
\label{sec:c5:adaptations}

\figref{fig:c5:adaptation-motivation} summarises ambiguities motivating the controlled adaptations. These methods change the Stage~2 input or supervision while retaining the principal Direct SFT optimisation configuration.

\begin{figure}[htb!]
    \centering
    \includegraphics[width=0.95\linewidth]{figures/adaptation_motivation.pdf}
    \caption{Representative challenges motivating the Stage~2 adaptations: fine-grained anatomical ambiguity, temporal evidence absent from an isolated frame, plausible outputs outside the benchmark ontology, and cases in which the available image evidence is insufficient to distinguish neighbouring structures.}
    \label{fig:c5:adaptation-motivation}
\end{figure}

\subsubsection{Sequential Dialogue Prediction}

Sequential Dialogue Prediction factorises the structured response into two turns, following the motivation of divide-and-discern dialogue tuning~\cite{zhang2026bootstrapping}. The first turn predicts one verb for each grounded instance. The second predicts the targets while retaining the first response as dialogue context. During training, target prediction is conditioned on the ground-truth verbs and supervised loss is applied to both assistant responses. During inference, the second turn receives the verbs generated in the first turn, so verb errors can propagate to target prediction. The model is trained independently from the same public base model as Direct SFT, with otherwise identical QLoRA and optimisation settings.

\begin{figure}[htb!]
    \centering
    \includegraphics[width=0.95\linewidth]{figures/sequential_dialogue_prediction.pdf}
    \caption{Direct interaction prediction and Sequential Dialogue Prediction. Direct SFT predicts verbs and targets jointly. The sequential formulation first predicts verbs and then predicts targets conditioned on the preceding dialogue response.}
    \label{fig:c5:sequential-dialogue}
\end{figure}

\subsubsection{Temporal Visual Context}

The temporal variants retain centre-frame interaction prediction but replace the single image with an ordered native multimodal input. The three-frame model uses $(t-1,t,t+1)$ and the five-frame model uses $(t-2,t-1,t,t+1,t+2)$, where neighbouring inputs are immediately adjacent source-video frames. Every frame contains overlays derived from its own instrument masks; no temporal tracking or mask propagation is used. The prompt explicitly identifies the current centre frame as the prediction target and states that the neighbouring frames provide context only. Consequently, only the centre frame supplies triplet supervision. Windows remain within the same video and split, and at video boundaries only available neighbouring frames are used.

The total visual-token budget is held equal to the single-frame setting and is therefore shared across the temporal frames. The comparison between three and five frames is consequently exploratory rather than a controlled temporal-window ablation: the five-frame input contains more visual content within the same budget. Apart from the temporal input and validation schedule, training follows the Direct SFT configuration. An earlier cached-feature implementation was technically inappropriate and is not included in the study.

\subsubsection{Anatomical Context Priors}

This adaptation supplies coarse anatomy around each instrument tip as text in addition to the standard visual input. It uses the same EndoViT checkpoint and label space described in Chapter~\ref{chap:triplet-grounding}. A heuristic estimates the distal endpoint from the instrument mask, entry direction, shaft orientation, and laparoscopic field boundary. A circular neighbourhood centred on the endpoint has radius one-eighth of the smaller image dimension. Within this region, each pixel is assigned its highest-probability EndoViT class; classes are ranked by their proportion of the region, and the three most prevalent classes and proportions are appended to the prompt. Anatomical context is supplied during both training and inference. The model is trained independently from the same base model with the Direct SFT configuration.

\begin{figure}[htb!]
    \centering
    \includegraphics[width=0.98\linewidth]{figures/instrument_tip_prior-examples.pdf}
    \caption{Representative examples of the heuristic instrument-tip estimation used to extract local anatomical context. The estimated instrument tip (red dot) defines the centre of a circular region (yellow) from which the local anatomical distribution is obtained using EndoViT predictions. The accompanying text reports the three most prevalent predicted anatomical classes and their pixel proportions within the selected region. (A) The heuristic estimates the tips of both visible instruments, producing anatomically relevant local context. (B) A limitation of the heuristic: the grasper contains two plausible distal extremities, making the true contact point ambiguous. The selected tip may consequently yield less representative anatomical context. Qualitative inspection nevertheless indicated that the heuristic generally produced reasonable tip estimates across the examined examples.}
    \label{fig:c5:instrument-tip-prior}
\end{figure}

\subsubsection{Ontology-Guided Prompting}

Ontology-Guided Prompting augments the instruction with the valid verb--target combinations for each supplied instrument class. These combinations are taken without modification from the 100-triplet ontology introduced in Chapter~\ref{chap:triplet-grounding}. Ontology information is included during both fine-tuning and evaluation, but decoding remains autoregressive and unconstrained. The experiment therefore evaluates prompt-level access to the ontology rather than constrained prediction.

\subsubsection{Distilled-Reasoning Supervision}

Qwen3.7-Max, a closed-weight commercial teacher model, was prompted with the grounded image and ground-truth labels to generate a structured rationale for each instance. Each response describes contact evidence, action evidence, a primary target hypothesis, a contrastive target hypothesis, uncertainty, and the final triplet. Rationales were generated for all splits, but only training-split rationales were used as supervision. Approximately ten rationales were inspected informally; they appeared plausible, but neither systematic quality control nor clinical expert validation was performed.

\begin{figure}[htb!]
    \centering
    \includegraphics[width=0.95\linewidth]{figures/distilled_reasoning_supervision.pdf}
    \caption{Distilled-reasoning supervision. Qwen3.7-Max generates structured rationales from grounded images and ground-truth triplets. A Qwen3-VL-8B-Instruct student is then trained to reproduce the rationale and final structured prediction.}
    \label{fig:c5:distilled-reasoning}
\end{figure}

The Distilled Rationale Only model is trained exclusively to generate the complete rationale followed by the final structured prediction. The multitask model uses equal numbers of Direct SFT and distilled-rationale examples. The same trained multitask model is evaluated in two modes: \emph{Multitask---Direct SFT} requests the concise response, whereas \emph{Multitask---Distilled Rationale} requests the rationale before the final prediction. Only the final structured triplet is scored. Rationale and answer tokens have equal loss weight, so the longer rationale accounts for most supervised response tokens. All student models otherwise follow the Direct SFT configuration.

\section{Experimental Protocol}
\label{sec:c5:experimental-protocol}

\subsection{Dataset and Samples}

All experiments use the official CholecTriplet-Seg training, validation, and test partitions introduced in Chapter~\ref{chap:triplet-grounding}: 22,818 training frames from 45 videos, 2,496 validation frames from one video, and 5,641 test frames from four videos. The chapter reuses the established annotations without further relabelling or cleaning. Frames containing no visible instrument instance are excluded because Stage~2 predicts interactions only for grounded instances. Each single-frame sample contains all grounded instances and associated triplets in one annotated frame. All adaptations use the same data partitions as Direct SFT.

\subsection{Evaluation}

Performance is measured using the mask-based IVTMetrics segmentation protocol defined in Chapter~\ref{chap:triplet-grounding}. The reported measures are instrument AP (AP$_I$), verb AP (AP$_V$), target AP (AP$_T$), instrument--verb AP (AP$_{IV}$), instrument--target AP (AP$_{IT}$), and complete-triplet AP (AP$_{IVT}$). The same matching, overlap, assignment, and class-averaging rules are used here.

Ground-truth grounding fixes the instrument masks, classes, and instance count; AP$_I$ is therefore 100 by construction, while the remaining metrics are computed from Stage~2 outputs. Predicted grounding uses the fixed Stage~1 detector and includes both grounding and interaction errors. Only Direct SFT is evaluated with predicted instances; all adaptations are evaluated under ground-truth grounding. Every reported experiment is a single run with seed 42. No confidence intervals, repeated-seed estimates, or significance tests are available, so small differences are interpreted as observations rather than robust improvements.

\section{Results}
\label{sec:c5:results}

\subsection{Interaction Prediction with Ground-Truth Grounding}

\tabref{tab:c5:principal-adaptations} reports the principal comparison under ground-truth grounding. Direct SFT reaches AP$_{IVT}=14.56$, establishing that the MLLM can perform structured interaction prediction when instrument instances are supplied. Three-frame temporal context gives the highest observed AP$_{IVT}$ at 15.89, an absolute increase of 1.33 points, and also gives the highest AP$_V$ and AP$_{IV}$. Sequential Dialogue Prediction reaches AP$_{IVT}=15.31$, 0.75 points above Direct SFT, despite slightly lower component metrics. These are promising single-seed observations, not evidence of statistically reliable superiority.

\begin{table}[htb!]
\centering
\caption{Principal Stage~2 results under ground-truth instrument grounding. AP$_I=100$ by construction and is omitted. Bold values indicate the highest observed result in each column.}
\label{tab:c5:principal-adaptations}
\begin{tabular}{lccccc}
\toprule
Method & AP$_V$ & AP$_T$ & AP$_{IV}$ & AP$_{IT}$ & AP$_{IVT}$ \\
\midrule
Direct SFT & 55.99 & \textbf{26.35} & 31.03 & 21.87 & 14.56 \\
Sequential Dialogue Prediction & 55.92 & 26.21 & 30.48 & 21.61 & 15.31 \\
Anatomical Context Priors & 53.35 & 26.27 & 31.48 & \textbf{23.15} & 14.96 \\
Temporal Visual Context, 3 frames & \textbf{56.31} & 26.32 & \textbf{32.37} & 21.66 & \textbf{15.89} \\
\bottomrule
\end{tabular}
\end{table}

The anatomical-prior model produces a mixed result. It attains the highest AP$_{IT}$ in the principal table and AP$_{IVT}=14.96$, but AP$_T$ remains effectively unchanged from Direct SFT and AP$_V$ declines. The experiment therefore does not show that this prompt-level anatomical representation improves target classification. It instead suggests that coarse context and heuristic tip extraction can alter instance--target associations without providing a consistent overall benefit.

\subsection{Exploratory and Distilled-Supervision Results}

\tabref{tab:c5:secondary-adaptations} collects adaptations that yielded negative or mixed results. Five-frame context falls to AP$_{IVT}=11.61$. Because the same total visual budget is shared across more frames, this result shows that adding frames under this formulation does not automatically help; it does not isolate temporal-window length as a causal factor. Ontology-Guided Prompting reaches AP$_{IVT}=14.44$, providing no evidence that listing valid combinations in the prompt improves the principal outcome.

\begin{table}[htb!]
\centering
\caption{Exploratory adaptations under ground-truth grounding. The Direct SFT row is repeated as a reference. The two multitask rows evaluate the same trained model with different inference instructions.}
\label{tab:c5:secondary-adaptations}
\begin{tabular}{lccccc}
\toprule
Method & AP$_V$ & AP$_T$ & AP$_{IV}$ & AP$_{IT}$ & AP$_{IVT}$ \\
\midrule
Direct SFT & 55.99 & 26.35 & 31.03 & 21.87 & 14.56 \\
Temporal Visual Context, 5 frames & 53.52 & 22.10 & 27.72 & 17.39 & 11.61 \\
Ontology-Guided Prompting & 55.81 & 24.48 & 33.07 & 20.68 & 14.44 \\
Distilled Rationale Only & 53.22 & 20.88 & 27.30 & 13.59 & 8.69 \\
Multitask---Direct SFT & 52.86 & 25.68 & 31.70 & 21.60 & 14.57 \\
Multitask---Distilled Rationale & 54.90 & 25.18 & 33.20 & 22.49 & 14.32 \\
\bottomrule
\end{tabular}
\end{table}

Distilled Rationale Only supervision performs poorly, particularly on target-related measures. The multitask model recovers most Direct SFT performance, but neither evaluation mode improves AP$_{IVT}$ meaningfully: Direct SFT evaluation gives 14.57 and distilled-rationale evaluation gives 14.32. Thus, the experiments do not establish that the generated rationales are beneficial. One plausible explanation is that equal token weighting causes the long rationale to dominate the objective while the short final answer determines evaluation; this remains a hypothesis because answer-token reweighting was not tested.

\subsection{Complete Two-Stage System}

\tabref{tab:c5:predicted-grounding} evaluates Direct SFT with Stage~1 predictions and compares it with the Chapter~\ref{chap:triplet-grounding} systems under the same dataset and IVTMetrics protocol. The published comparison values are reproduced from that chapter. The two-stage system obtains AP$_I=80.64$, which is fixed by Stage~1 and is not attributable to the MLLM. It also produces substantially higher AP$_V$, AP$_T$, AP$_{IV}$, and AP$_{IT}$ than TargetFusionNet. Complete-triplet performance is essentially comparable: AP$_{IVT}=13.52$ for the two-stage system and 13.45 for TargetFusionNet. The 0.07-point difference is too small to support an improvement claim.

\begin{table}[htb!]
\centering
\caption{Mask-based triplet segmentation using predicted instrument instances. Comparison results are reproduced from Chapter~\ref{chap:triplet-grounding}. Bold values indicate the highest observed score in each column.}
\label{tab:c5:predicted-grounding}
\begin{tabular}{lcccccc}
\toprule
Method & AP$_I$ & AP$_V$ & AP$_T$ & AP$_{IV}$ & AP$_{IT}$ & AP$_{IVT}$ \\
\midrule
RDV-Det~\cite{nwoye2023cholectriplet2022} & 0.09 & 0.11 & 0.08 & 0.08 & 0.04 & 0.03 \\
Mask2Former-Triplet~\cite{alabi2026grounding} & 65.24 & 45.61 & 20.75 & 23.03 & 16.47 & 12.23 \\
TargetFusionNet~\cite{alabi2026grounding} & 67.19 & 46.23 & 23.70 & 24.82 & 17.74 & 13.45 \\
Two-stage Direct SFT & \textbf{80.64} & \textbf{51.89} & \textbf{25.27} & \textbf{29.01} & \textbf{18.73} & \textbf{13.52} \\
\bottomrule
\end{tabular}
\end{table}

The reduction from AP$_{IVT}=14.56$ under ground-truth grounding to 13.52 in the complete system reflects compounded errors. False-negative Stage~1 instances never reach Stage~2, false positives generate additional interaction predictions, and incorrect instrument identities cannot be corrected by the MLLM. Moreover, improvements in individual components need not translate proportionally to AP$_{IVT}$ because all three labels and the spatial association must be correct simultaneously.

\subsection{Sensitivity to Optimisation}

Exploratory Direct SFT runs varied training duration, effective batch size, and learning rate simultaneously. Across these configurations, AP$_{IVT}$ ranged from 14.56 to 16.05, with different settings favouring different component metrics. This demonstrates material sensitivity to the overall training configuration but cannot identify the causal effect of any individual hyperparameter. The later runs therefore are not substituted for the principal Direct SFT baseline or compared with adaptations as though they formed a controlled study.

\section{Discussion}
\label{sec:c5:discussion}

\subsection{Viability of the Two-Stage Formulation}

The principal finding is that an MLLM can act as a viable second-stage interaction predictor when instrument instances are supplied. Direct SFT requires no explicit rationale, anatomical prompt, or ontology constraint, yet provides a strong reference point under ground-truth grounding and remains competitive in the complete predicted-grounding pipeline. The decomposition focuses the MLLM on semantic interaction prediction while leaving dense localisation to a specialist model. This does not prove that two-stage modelling is universally preferable to end-to-end learning, but it establishes the formulation as credible and exposes the contribution and failure modes of each stage more clearly.

The complete-system results also clarify the relationship to TargetFusionNet. Component-level scores are substantially stronger for the two-stage Direct SFT pipeline, but AP$_{IVT}$ is essentially unchanged. Complete-triplet evaluation is conjunctive: better verb or target recognition contributes only when localisation, instrument identity, and the remaining semantic component are also correct for the same instance. Consequently, component gains should not be presented as equivalent to complete interaction gains.

\subsection{Task Factorisation and Temporal Evidence}

Sequential Dialogue Prediction and three-frame context produce the strongest observed AP$_{IVT}$ results among the principal Stage~2 experiments. Their effects are nevertheless metric-specific. Sequential prediction increases AP$_{IVT}$ while slightly decreasing every reported component or pair metric relative to Direct SFT. It may encourage more coherent complete predictions, but the present experiment does not quantify propagation from an incorrect first-turn verb, and the training--inference difference between ground-truth and predicted verb conditioning remains important.

Three-frame input gives a larger AP$_{IVT}$ increase and modest improvements in AP$_V$ and AP$_{IV}$. This is consistent with the premise that neighbouring frames can reveal motion or interaction progression absent from a still image. However, the experiment uses ground-truth grounding and only one seed. The five-frame result further shows that additional visual context is not automatically useful when a fixed token budget is divided among more frames. It should not be interpreted as a controlled demonstration that three frames are intrinsically optimal.

\subsection{Why Additional Priors and Longer Supervision Did Not Consistently Help}

The anatomical and ontology experiments provide useful negative evidence about prompt-level integration. EndoViT context is conceptually connected to the weak anatomical priors investigated in Chapter~\ref{chap:triplet-grounding}, but the present conversion of dense predictions into three textual class proportions does not improve AP$_T$. This result does not show that anatomy is unhelpful; it shows that this heuristic tip region and prompt representation do not reliably improve the final task.

Similarly, access to valid verb--target combinations does not improve AP$_{IVT}$. Because generation remains unconstrained, the model may not consistently use the supplied ontology. The experiment therefore evaluates contextual guidance, not the potential value of constrained decoding.

Distilled-reasoning supervision also fails to outperform Direct SFT. The rationales were generated at scale but were not clinically validated, and only a small informal sample was inspected. Equal weighting also directs most response-token loss towards reproducing the rationale rather than the evaluated triplet. These design properties provide plausible explanations for the negative result, but their individual effects were not tested and should not be treated as established causes.

\section{Limitations}
\label{sec:c5:limitations}

The study has six principal limitations. First, every result is based on one run with seed 42; small differences may reflect run-to-run variation, and no claim of statistical significance or robustness can be made. Second, only Direct SFT is evaluated with predicted grounding. The sequential, temporal, anatomical, ontology, and distilled-supervision variants are assessed only with ground-truth instances, so their behaviour in the complete system is unknown.

Third, adaptations are evaluated independently. No controlled combinations test whether sequential factorisation and temporal context are complementary. Fourth, optimisation sensitivity is substantial, and the exploratory optimisation runs change several variables at once. This complicates comparisons across adaptations developed under a shared but not exhaustively optimised configuration.

Fifth, malformed generations and verb-to-target error propagation were not quantified systematically. Invalid outputs receive no credit, but their frequency may differ between formulations and remains an unmeasured source of variability. Finally, all experiments use CholecTriplet-Seg and laparoscopic cholecystectomy. Generalisation to other procedures, acquisition conditions, and interaction ontologies has not been demonstrated.

\section{Future Research}
\label{sec:c5:future-work}

The immediate priority is repeated-seed evaluation of Direct SFT, Sequential Dialogue Prediction, and three-frame context, followed by evaluation of the strongest variants with predicted grounding. Temporal predicted grounding may require instance association or mask propagation because independently detected instances are not guaranteed to correspond across frames. More explicit temporal integration may also use motion and procedural continuity more effectively than distributing a fixed visual budget across independent images.

Sequential prediction should be tested with imperfect verb conditioning during training so that robustness to first-turn errors can be measured directly. Controlled combinations of the strongest adaptations would determine whether task factorisation and temporal evidence are complementary. Curriculum-style training could also initialise specialised variants from a converged Direct SFT adapter rather than training each independently from the base model.

For distilled supervision, future work should prioritise rationale quality control and objectives that increase the relative weight of final answer tokens. Comparison-based or contrastive auxiliary tasks may offer a more directly verifiable form of supervision than long free-text rationales. Finally, evaluation on additional procedures and datasets is required to establish whether two-stage interaction prediction transfers beyond cholecystectomy.

\section{Conclusion}
\label{sec:c5:conclusion}

This chapter formulates surgical triplet segmentation as specialist instrument grounding followed by MLLM-based interaction prediction. Direct SFT establishes the principal viable Stage~2 baseline. Under ground-truth grounding, Sequential Dialogue Prediction and three-frame temporal context produce promising observed gains in AP$_{IVT}$, although their component-level effects are mixed and their superiority is not established across seeds. Prompt-level anatomical priors, ontology guidance, five-frame context, and distilled-reasoning supervision do not improve performance consistently, demonstrating that more context or longer supervision is not inherently beneficial.

When the full system uses predicted masks, grounding errors reduce performance. Nevertheless, the two-stage Direct SFT pipeline produces substantially stronger component scores and complete-triplet AP comparable to TargetFusionNet. The results therefore support two-stage multimodal interaction prediction as a defensible culmination of the thesis: Chapters~\ref{chap:cholecinstanceseg} and~\ref{chap:triplet-grounding} establish the spatial substrate, and this chapter shows how semantic interaction prediction can be studied as a specialised second stage. The remaining limitations define a focused continuation beyond the thesis without making the present contribution dependent on unperformed experiments.
